\chapter{Genetic Algorithm Implementation}
\label{ch:ga}

A \ac{GA}~\cite{holland1975,goldberg1989} evolves a \emph{population} of candidate designs instead
of a single one. In every generation the designs are scored, the better ones are selected as
parents, their genes are recombined (crossover) and slightly perturbed (mutation), and the children
form the next generation. The \ac{GA} plugs into the same Algorithm $\rightarrow$ Feasibility
$\rightarrow$ Engine pipeline as \ac{SA} (Figure~\ref{fig:arch}): it only proposes decision vectors
and reads back their cost, under the same evaluation budget $B$, so the two algorithms are compared
on equal terms. Like \ac{SA}, it is written from scratch with NumPy only.

%=====================================================================================
\section{Chromosome representation}
\label{sec:ga-chromosome}

The \ac{GA} uses the decision vector $\mathbf{x}$ of Eq.~(\ref{eq:x}) directly as its chromosome;
no re-encoding is needed. Four levels are distinguished (Figure~\ref{fig:chromosome}):

\begin{itemize}
  \item \textbf{Gene:} one entry of $\mathbf{x}$, e.g.\ the orbit radius $a_2$ of slot~2.
  \item \textbf{Slot} $\mathbf{y}_i=[a_i,\iota_i,\Omega_i,M_{0,i},s_i]$: the five genes of one
        spacecraft (its circular orbit and its on/off switch).
  \item \textbf{Chromosome} $\mathbf{x}=[T,\mathbf{y}_1,\dots,\mathbf{y}_{N_{\max}}]$: one
        \emph{individual}, i.e.\ one complete mission design (the whole fleet and the mission
        duration), with $\dim\mathbf{x}=1+5N_{\max}$ genes (26 for $N_{\max}=5$, 51 for $N_{\max}=10$).
  \item \textbf{Population:} the $P$ chromosomes (candidate missions) of one generation.
\end{itemize}

\paragraph{Why a chromosome is a whole mission and not one spacecraft.} The fitness $J$
(Eq.~\ref{eq:J}) is only defined for a complete design: the coverage $C$ is the area imaged by the
\emph{union} of all spacecraft (two spacecraft imaging the same region are worth less than two
imaging different regions), the fleet cost $Z$ depends on how many spacecraft fly, the separation
constraint g3 couples \emph{pairs} of spacecraft, and the duration $T$ is shared by the whole fleet.
The unit that is evaluated and selected must therefore be the full mission.

\paragraph{Why a fixed length.} The fleet size is not stored; it follows from the switches,
$N=\sum_i s_i$ (Eq.~\ref{eq:N}). Every chromosome has the same $1+5N_{\max}$ genes whatever its
$N$, so standard crossover and mutation apply between any two parents, including parents with
different fleet sizes, and the \ac{GA} searches exactly the same space as \ac{SA}.

\paragraph{Why real coding.} The genes are kept as real numbers (and a binary switch) instead of
being converted to bit strings: a binary code would add Hamming cliffs (neighbouring values that
differ in many bits), a fixed precision, and would break the periodicity of the angles $\Omega_i$
and $M_{0,i}$. The existing \texttt{repair} operator already handles every gene kind: $T$, $a_i$
and $\iota_i$ are reflected at their bounds, $\Omega_i$ and $M_{0,i}$ wrap at $360^\circ$, and the
switches are rounded to $0/1$ with at least one switched on.

\paragraph{Active and dormant genes.} The orbit genes of a switched-off slot are carried along
but do not affect $J$ (\emph{dormant} genes). At any time the \emph{active} genes are $T$, all
$N_{\max}$ switches and the four orbit genes of each flying spacecraft, $1+N_{\max}+4N$ in total
(18 of the 26 genes in Figure~\ref{fig:chromosome}). Mutation acts on active genes only, so no
evaluation is spent on a change that cannot alter the fitness; the dormant genes are still
inherited, so switching a slot back on can revive an orbit found earlier in the search.

\begin{figure}[htbp]
\centering
\resizebox{\textwidth}{!}{%
\begin{tikzpicture}[
    font=\sffamily,
    swatch/.style={minimum width=0.38cm, minimum height=0.26cm, inner sep=0pt, rounded corners=1pt},
    legend/.style={font=\sffamily\scriptsize, inner sep=1pt, anchor=west},
    brace/.style={decorate, decoration={brace, mirror, amplitude=4pt}, line width=0.7pt, gray!70!black},
  ]
  % ---------------------------------------------------------------- (a) one chromosome
  \node[font=\sffamily\bfseries\small, anchor=west] at (-0.1,2.05) {(a) One chromosome = one individual = one complete mission design};

  % head gene: mission duration T
  \fill[headc, rounded corners=3pt] (0,0) rectangle (1.0,0.95);
  \node[text=white, font=\small] at (0.5,0.66) {$T$};
  \node[text=white, font=\scriptsize] at (0.5,0.27) {1.4\,d};
  \node[font=\sffamily\scriptsize\bfseries, text=headc] at (0.5,1.32) {Mission};
  \node[font=\sffamily\tiny, text=headc] at (0.5,1.08) {duration};

  % five slots: id / colour / switch / a / inc / raan / M0
  \foreach \sid/\col/\sw/\va/\vi/\vO/\vM in {1/sc1/1/50/30/0/0, 2/sc2/1/45/90/60/120,
                                             3/offc/0/63/45/300/10, 4/sc4/1/40/120/200/240,
                                             5/offc/0/80/10/90/45} {
    \pgfmathsetmacro{\x}{1.22+3.32*(\sid-1)}
    \ifnum\sw=1
      \def\tc{black}\def\swc{white}\def\stxt{flying}\tikzset{bord/.style={solid}}
    \else
      \def\tc{gray}\def\swc{gray!60!black}\def\stxt{switched off: dormant genes}\tikzset{bord/.style={dashed}}
    \fi
    % gene kinds: bounded reals (a, inc) | angles (raan, M0) | switch (s)
    \fill[fill=\col!22] (\x,0) rectangle ++(1.24,0.95);
    \fill[fill=\col!45] ({\x+1.24},0) rectangle ++(1.24,0.95);
    \fill[fill=\col!85] ({\x+2.48},0) rectangle ++(0.62,0.95);
    \foreach \k in {1,...,4} { \draw[white, line width=1pt] ({\x+0.62*\k},0) -- ++(0,0.95); }
    \draw[draw=\col!85!black, line width=0.9pt, rounded corners=2pt, bord] (\x,0) rectangle ++(3.1,0.95);
    % gene symbols (top) and example values (bottom)
    \node[font=\small, text=\tc] at ({\x+0.31},0.66) {$a_{\sid}$};
    \node[font=\small, text=\tc] at ({\x+0.93},0.66) {$\iota_{\sid}$};
    \node[font=\small, text=\tc] at ({\x+1.55},0.66) {$\Omega_{\sid}$};
    \node[font=\small, text=\tc] at ({\x+2.17},0.66) {$M_{0,\sid}$};
    \node[font=\small, text=\swc] at ({\x+2.79},0.66) {$s_{\sid}$};
    \node[font=\scriptsize, text=\tc] at ({\x+0.31},0.27) {\va};
    \node[font=\scriptsize, text=\tc] at ({\x+0.93},0.27) {$\vi^\circ$};
    \node[font=\scriptsize, text=\tc] at ({\x+1.55},0.27) {$\vO^\circ$};
    \node[font=\scriptsize, text=\tc] at ({\x+2.17},0.27) {$\vM^\circ$};
    \node[font=\scriptsize\bfseries, text=\swc] at ({\x+2.79},0.27) {\sw};
    % header
    \node[font=\sffamily\scriptsize\bfseries, text=\col!85!black] at ({\x+1.55},1.32) {Slot \sid{} $\cdot$ spacecraft \sid};
    \node[font=\sffamily\tiny, text=\col!85!black] at ({\x+1.55},1.08) {\stxt};
  }

  % braces: gene, slot, chromosome
  \draw[brace] (1.24,-0.08) -- (1.82,-0.08) node[midway, below=4pt, font=\sffamily\scriptsize, align=center] {one gene\\[-1pt]($a_1$ = 50\,km)};
  \draw[brace] (4.56,-0.08) -- (7.62,-0.08) node[midway, below=4pt, font=\sffamily\scriptsize, align=center] {one slot = one spacecraft\\[-1pt]$\mathbf{y}_2$: 5 genes};
  \node[font=\sffamily\scriptsize, anchor=west, align=left] at (9.2,-0.45)
       {$N=\sum_i s_i = 1+1+0+1+0 = 3$ spacecraft fly;\\[-1pt]
        active genes: $1+N_{\max}+4N = 18$ of 26};
  \draw[brace, line width=0.9pt, decoration={brace, mirror, amplitude=6pt}] (0,-1.15) -- (17.58,-1.15)
       node[midway, below=7pt, font=\sffamily\small] {chromosome $\mathbf{x}$: $\dim\mathbf{x} = 1 + 5N_{\max} = 1 + 5\times5 = 26$ genes};

  % legend
  \node[swatch, fill=headc] (l1) at (0.25,-2.3) {};
  \node[legend, right=2pt of l1] (t1) {mission duration (bounded real)};
  \node[swatch, fill=sc1!22, right=12pt of t1] (l2) {};
  \node[legend, right=2pt of l2] (t2) {bounded real, reflected at the bounds ($a_i$, $\iota_i$)};
  \node[swatch, fill=sc1!45] (l3) at (0.25,-2.75) {};
  \node[legend, right=2pt of l3] (t3) {angle, wraps at $360^\circ$ ($\Omega_i$, $M_{0,i}$)};
  \node[swatch, fill=sc1!85, right=12pt of t3] (l4) {};
  \node[legend, right=2pt of l4] (t4) {binary switch ($s_i$)};
  \node[swatch, fill=offc!22, draw=offc!85!black, dashed, right=12pt of t4] (l5) {};
  \node[legend, right=2pt of l5] {switched-off slot};

  % ---------------------------------------------------------------- (b) the population
  \node[font=\sffamily\bfseries\small, anchor=west] at (-0.1,-3.55) {(b) Population = $P$ chromosomes (candidate missions) per generation};
  \def\xo{1.1}
  \foreach \r/\lab/\bA/\bB/\bC/\bD/\bE/\n in {0/1/1/1/0/1/0/3, 1/2/0/1/1/0/0/2, 2/3/1/1/1/1/1/5,
                                               3/4/1/0/0/0/0/1, 4/5/0/1/0/1/1/3, 6/P/0/0/1/1/0/2} {
    \pgfmathsetmacro{\y}{-4.35-0.45*\r}
    \fill[fill=headc, rounded corners=1pt] (\xo,\y) rectangle ++(0.3,0.3);
    \foreach \j/\bit in {1/\bA, 2/\bB, 3/\bC, 4/\bD, 5/\bE} {
      \pgfmathsetmacro{\xs}{\xo+0.4+0.9*(\j-1)}
      \ifnum\bit=1 \def\pc{sc1}\else\def\pc{offc}\fi
      \fill[fill=\pc!22] (\xs,\y) rectangle ++(0.32,0.3);
      \fill[fill=\pc!45] ({\xs+0.32},\y) rectangle ++(0.32,0.3);
      \fill[fill=\pc!85] ({\xs+0.64},\y) rectangle ++(0.16,0.3);
    }
    \node[font=\sffamily\scriptsize, anchor=east] at ({\xo-0.08},{\y+0.15}) {\lab};
    \node[font=\sffamily\scriptsize, anchor=west] at ({\xo+4.95},{\y+0.15}) {$N=\n$};
  }
  \node[font=\small] at ({\xo+2.4},-6.45) {$\vdots$};
  \draw[draw=sc1!70!black, line width=0.9pt, rounded corners=2pt] ({\xo-0.06},-4.41) rectangle ({\xo+5.85},-3.99);
  \draw[decorate, decoration={brace, amplitude=5pt}, line width=0.7pt, gray!70!black] (0.45,-7.05) -- (0.45,-4.0)
       node[midway, left=6pt, font=\sffamily\small, text=black] {$P$};
  \node[font=\sffamily\scriptsize, text=sc1!70!black, anchor=west] at ({\xo+6.0},-4.2) {$\leftarrow$ the chromosome of (a)};

  \node[font=\sffamily\scriptsize, anchor=north west, text width=9.2cm, align=left] at (8.3,-4.6)
     {Each row is one candidate fleet. Individuals differ in their orbits \emph{and} in their
      fleet size $N$, because $N$ is set by the switches. In every generation all $P$ individuals
      are evaluated by the pipeline (Feasibility $\rightarrow$ Engine); the better ones are
      selected as parents, and crossover and mutation produce the next generation.};
\end{tikzpicture}%
}
\caption{Chromosome representation of the \ac{GA}. (a)~A gene is one number, a slot is the five
genes of one spacecraft, and the chromosome is the whole decision vector (Eq.~\ref{eq:x}): the
mission duration followed by $N_{\max}$ slots. In this example ($N_{\max}=5$) slots 1, 2 and 4 fly
($N=3$); the orbit genes of slots 3 and 5 are dormant but still inherited. (b)~The population holds
$P$ such chromosomes, each a complete mission with its own fleet size.}
\label{fig:chromosome}
\end{figure}

%=====================================================================================
\section{Population size}
\label{sec:ga-popsize}

The population size $P$ is set from the length of the chromosome:
\begin{equation}
P=\min\!\Bigl(60,\;\max\!\Bigl(30,\;10\Bigl\lceil\tfrac{1.5\,\dim\mathbf{x}}{10}\Bigr\rceil\Bigr)\Bigr)
\quad\Rightarrow\quad
P=40\ \ (N_{\max}=5,\ \dim\mathbf{x}=26),\qquad P=60\ \ (N_{\max}=10,\ \dim\mathbf{x}=51).
\label{eq:popsize}
\end{equation}
With elitism (the $n_e$ best individuals are copied unchanged into the next generation), each
generation after the first costs $P-n_e$ evaluations, so a budget of $B$ evaluations gives
\begin{equation}
G = 1+\Bigl\lceil\frac{B-P}{P-n_e}\Bigr\rceil
\label{eq:generations}
\end{equation}
generations. Table~\ref{tab:popsize} lists the resulting settings. The choice rests on four
arguments.

\begin{table}[htbp]
\centering
\caption{Population size and number of generations for the budget $B=2000$ used by \ac{SA}, with
the default elite share $\varepsilon=5\,\%$ (Section~\ref{sec:ga-split}).}
\label{tab:popsize}
\begin{tabular}{@{}lccccc@{}}
\toprule
Case studies & $N_{\max}$ & $\dim\mathbf{x}$ & $P$ & $n_e$ & $G$ \\
\midrule
Eros 10\,d / 1\,d, Kleopatra 3\,d                           & 5  & 26 & 40 & 2 & 53 \\
Eros 1\,d / 12\,h / 6\,h / 2\,h, Kleopatra 3\,d / 12\,h ($N_{\max}=10$) & 10 & 51 & 60 & 3 & 36 \\
\bottomrule
\end{tabular}
\end{table}

\paragraph{1. Diversity across the genes.} A population that is small compared with the number
of genes cannot hold enough different values of every gene at once; crossover then has little to
recombine and the population collapses onto one design early (premature convergence). A common
rule of thumb is a population between one and two times the number of decision variables;
$P\approx1.5\dim\mathbf{x}$ sits in the middle of that range.

\paragraph{2. Every fleet size must be represented.} The fleet size is encoded by the switches.
In the random initial population each switch is on with probability $\tfrac12$, so $N$ follows a
binomial distribution (an all-off draw is repaired to $N=1$). Table~\ref{tab:popN} gives the
expected number of individuals per fleet size. With $P=40$ and $N_{\max}=5$ every fleet size from
1 to 4 starts with at least six individuals, so the competition between fleet sizes is decided by
selection and not by chance. For $N_{\max}=10$ small fleets are rare in a uniform start
(fewer than one individual with $N=1$, about three with $N=2$), although \ac{SA} found that two to
four spacecraft are best (Chapter~\ref{ch:results}); the initial population is therefore
spread evenly over the fleet sizes (Section~\ref{sec:ga-init}).

\begin{table}[htbp]
\centering
\caption{Expected number of individuals per fleet size $N$ in a uniformly random initial
population, $P\cdot\Pr(N)$ with $N\sim\mathrm{Binomial}(N_{\max},\tfrac12)$ and an all-off draw
counted as $N=1$.}
\label{tab:popN}
\footnotesize
\begin{tabular}{@{}lccccccccccc@{}}
\toprule
 & \multicolumn{10}{c}{fleet size $N$} \\
\cmidrule(l){2-11}
 & 1 & 2 & 3 & 4 & 5 & 6 & 7 & 8 & 9 & 10 \\
\midrule
$N_{\max}=5$, $P=40$  & 7.5 & 12.5 & 12.5 & 6.3 & 1.3 & -- & -- & -- & -- & -- \\
$N_{\max}=10$, $P=60$ & 0.6 & 2.6 & 7.0 & 12.3 & 14.8 & 12.3 & 7.0 & 2.6 & 0.6 & 0.1 \\
\bottomrule
\end{tabular}
\end{table}

\paragraph{3. A fixed budget forces a trade-off.} \ac{GA} and \ac{SA} receive the same budget
$B$, so $P\cdot G\approx B$: every individual added to the population removes evaluations from
later generations. A small population converges early; a large one leaves too few generations
for selection and crossover to act. In the limit $P=B$ the \ac{GA} is a single random generation,
i.e.\ random search, which validation V4 already showed to be clearly worse than \ac{SA}. The cap
of 60 for $N_{\max}=10$ follows from this: $1.5\times51\approx77$ individuals would leave only about
28 generations.

\paragraph{4. The cost of a generation.} The engine caches the simulation of every spacecraft by
its four orbit genes (Chapter~\ref{ch:sa}). A child that inherits whole slots from its parents
reuses their cached simulations, so the cost of an evaluation comes mainly from mutated orbits,
not from the population size. A population holds at most $P\,N_{\max}$ distinct orbits (200 for
$N_{\max}=5$, 600 for $N_{\max}=10$); the cache size is raised from 256 to 1024 orbits so that the
parents' simulations are not evicted while their children are evaluated. Caching does not change
any fitness value, only the run time.

The value of Eq.~(\ref{eq:popsize}) is a starting point. It is checked in the parameter study by
repeating the runs with $P=10$, 20, 40, 80 and 120 at the same budget.

%=====================================================================================
\section{Initial population}
\label{sec:ga-init}

Two ways of drawing the first generation are compared:
\begin{itemize}
  \item \textbf{Uniform:} every gene is drawn uniformly within its bounds and every switch is a
        coin flip. The fleet size then follows a binomial distribution (Table~\ref{tab:popN}):
        medium fleets are common, small and large fleets rare.
  \item \textbf{Stratified} (default): the population is first divided into $N_{\max}$ equal
        groups, one per fleet size $N=1,\dots,N_{\max}$ (the remainder of $P/N_{\max}$, if any, is
        assigned to randomly chosen sizes). In the group of size $N$, exactly $N$ randomly chosen
        slots are switched on; all orbit genes and $T$ are then drawn uniformly as before.
\end{itemize}

\begin{table}[htbp]
\centering
\caption{Number of individuals per fleet size $N$ in the first generation (uniform: expected value).}
\label{tab:init}
\footnotesize
\begin{tabular}{@{}llcccccccccc@{}}
\toprule
 & & \multicolumn{10}{c}{fleet size $N$} \\
\cmidrule(l){3-12}
 & & 1 & 2 & 3 & 4 & 5 & 6 & 7 & 8 & 9 & 10 \\
\midrule
\multirow{2}{*}{$N_{\max}=5$, $P=40$}  & uniform    & 7.5 & 12.5 & 12.5 & 6.3 & 1.3 & -- & -- & -- & -- & -- \\
                                       & stratified & 8 & 8 & 8 & 8 & 8 & -- & -- & -- & -- & -- \\
\midrule
\multirow{2}{*}{$N_{\max}=10$, $P=60$} & uniform    & 0.6 & 2.6 & 7.0 & 12.3 & 14.8 & 12.3 & 7.0 & 2.6 & 0.6 & 0.1 \\
                                       & stratified & 6 & 6 & 6 & 6 & 6 & 6 & 6 & 6 & 6 & 6 \\
\bottomrule
\end{tabular}
\end{table}

\paragraph{Why stratified.} The fleet size is the most important single decision of the design:
it sets the fleet cost $Z$ directly and the coverage through the number of cameras. A uniform start
already prefers some sizes before any design has been evaluated: with $N_{\max}=10$ it contains on
average 0.6 designs with one spacecraft and 2.6 with two, but 14.8 with five, although \ac{SA} found
two to four spacecraft to be best. The \ac{GA} would first have to reduce the fleet size step by step
through switch flips. A stratified start gives every fleet size the same number of candidates, so
that the comparison between fleet sizes is made by selection, on the fitness, and not by the
sampling. It costs nothing: the orbits are drawn exactly as before.

\paragraph{Feasibility of the start.} Every individual is drawn until it passes the pre-engine
checks (radius band, duration and separation), at most 200 times, as \ac{SA}'s warm start does. A
test on 200 random designs per fleet size showed that at least 93\,\% of them pass even with ten
spacecraft (Eros, $T_{\max}$ from 6\,h to 10\,d), so the start is not biased towards small fleets by
rejections.

Both start types are compared in the parameter study.

%=====================================================================================
\section{Selection}
\label{sec:ga-selection}

Parents are chosen by tournament selection: $k$ individuals are drawn at random from the whole
population (with replacement) and the best of them, compared feasible first and then by lower $J$,
becomes a parent. The default is the binary tournament, $k=2$, with the selection probabilities of
Eq.~(\ref{eq:tournament}). For crossover, the second parent is drawn again until it is a different
individual from the first, since crossing a design with itself only produces a copy.

Tournament selection is used instead of fitness-proportional (roulette-wheel) selection because of
the scale of the cost: a feasible design has $J\in[0,1]$, while a design that fails a constraint
costs at least $10^3$ (Eq.~\ref{eq:penalty}). Roulette selection needs positive fitness values on a
meaningful scale and would have to transform this cost; a tournament only compares two designs,
so it is unaffected by the scale and applies the feasibility-first rule of \ac{SA} exactly. Its
selection pressure is set by $k$ alone: a larger $k$ favours the best designs more strongly and
reduces the diversity faster; $k=2$, 3 and 5 are compared in the parameter study.

%=====================================================================================
\section{Population split: elites, crossover and mutation}
\label{sec:ga-split}

Each new generation is assembled from three groups (Figure~\ref{fig:split}):
\begin{equation}
n_e=\operatorname{round}(\varepsilon P),\qquad
n_c=\operatorname{round}\bigl(\chi\,(P-n_e)\bigr),\qquad
n_m=P-n_e-n_c,
\label{eq:split}
\end{equation}
where $\varepsilon$ is the \emph{elite share} and $\chi$ the \emph{crossover fraction}. The defaults
are $\varepsilon=5\,\%$ and $\chi=0.8$ (the same structure as the elite count and crossover
fraction of MATLAB's \texttt{ga}); Table~\ref{tab:split} gives the resulting group sizes.

\begin{itemize}
  \item \textbf{Elites} ($n_e$): the $n_e$ best \emph{distinct} individuals (ranked feasible first, then by
        lower $J$; an infeasible design costs more than any feasible one, Eq.~\ref{eq:penalty}) are copied unchanged. Their fitness is
        already known, so they are not evaluated again.
  \item \textbf{Crossover children} ($n_c$): two parents are chosen by tournament, their slots are
        aligned, and a child is assembled from the slots of both
        (Section~\ref{sec:ga-crossover}). If $n_c$ is odd, the last pair contributes one child.
        A child identical to one of its parents is mutated, so that no evaluation is spent on a copy.
  \item \textbf{Mutation children} ($n_m$): one parent is chosen by tournament, copied, and one or
        two of its active genes are changed: a Gaussian step on $T$, $a_i$, $\iota_i$, $\Omega_i$
        or $M_{0,i}$, or a flip of a switch $s_i$; the result is passed through \texttt{repair}.
\end{itemize}
Each generation therefore costs $n_c+n_m=P-n_e$ evaluations, as used in Eq.~(\ref{eq:generations}).

\begin{figure}[htbp]
\centering
\resizebox{\textwidth}{!}{%
\begin{tikzpicture}[
    font=\sffamily,
    op/.style={rounded corners=4pt, line width=0.9pt, text width=5.0cm, align=center,
               inner sep=5pt, font=\sffamily\scriptsize},
    arr/.style={-{Stealth[length=6pt]}, line width=0.9pt},
  ]
  % ---------------------------------------------------- generation g, ranked
  \foreach \k in {1,...,40} {
    \pgfmathtruncatemacro{\pct}{100-100*(\k-1)/39}
    \pgfmathsetmacro{\y}{8.0-0.2*\k}
    \fill[fill=sc2!\pct!offc!80] (0,\y) rectangle ++(1.0,0.2);
    \draw[white, line width=0.3pt] (0,\y) -- ++(1.0,0);
  }
  \draw[gray!70!black, line width=0.8pt, rounded corners=2pt] (0,0) rectangle (1.0,8.0);
  \draw[headc, line width=1.4pt, rounded corners=2pt] (-0.05,7.58) rectangle (1.05,8.05);
  \node[font=\sffamily\small\bfseries, align=center] at (0.5,8.75) {Generation $g$};
  \node[font=\sffamily\scriptsize, align=center] at (0.5,8.35) {$P=40$, ranked by $J$};
  \node[font=\sffamily\scriptsize, anchor=east] at (-0.1,7.8) {best};
  \node[font=\sffamily\scriptsize, anchor=east] at (-0.1,0.2) {worst};

  % ---------------------------------------------------- the three operators
  \node[op, draw=headc, fill=headc!8] (el) at (6.2,7.55)
       {\textbf{Elitism}\\ copy the $n_e$ best distinct designs unchanged\\ (no new evaluation)};
  \node[op, draw=sc1, fill=sc1!8] (cx) at (6.2,4.6)
       {\textbf{Crossover}\\ tournament $\rightarrow$ 2 parents $\rightarrow$ align slots
        $\rightarrow$ take each slot from parent A or B (or blend a matched pair)};
  \node[op, draw=sc4, fill=sc4!8] (mu) at (6.2,1.4)
       {\textbf{Mutation}\\ tournament $\rightarrow$ 1 parent $\rightarrow$ change 1--2 active genes
        (Gaussian step or switch flip) $\rightarrow$ repair};

  \draw[arr, headc] (1.05,7.8) -- (el.west);
  \draw[arr, sc1] (1.05,4.6) -- node[above, font=\sffamily\scriptsize, text=black] {binary tournament} (cx.west);
  \draw[arr, sc4] (1.05,1.4) -- node[above, font=\sffamily\scriptsize, text=black] {binary tournament} (mu.west);
  \node[font=\sffamily\scriptsize, text=gray!40!black, align=center] at (2.3,3.0) {all $P$ individuals\\ can be parents};

  % ---------------------------------------------------- generation g+1
  \fill[fill=headc] (10.4,7.6) rectangle (11.4,8.0);
  \fill[fill=sc1!70] (10.4,1.6) rectangle (11.4,7.6);
  \fill[fill=sc4!70] (10.4,0) rectangle (11.4,1.6);
  \foreach \k in {1,...,39} { \draw[white, line width=0.3pt] (10.4,{0.2*\k}) -- ++(1.0,0); }
  \draw[gray!70!black, line width=0.8pt, rounded corners=2pt] (10.4,0) rectangle (11.4,8.0);
  \node[font=\sffamily\small\bfseries] at (10.9,8.75) {Generation $g+1$};
  \node[font=\sffamily\scriptsize] at (10.9,8.35) {$P=40$};

  \draw[arr, headc] (el.east) -- (10.35,7.8);
  \draw[arr, sc1] (cx.east) -- (10.35,4.6);
  \draw[arr, sc4] (mu.east) -- (10.35,0.8);

  \node[font=\sffamily\scriptsize, anchor=west, align=left] at (11.55,7.8) {$n_e=2$ elites};
  \node[font=\sffamily\scriptsize, anchor=west, align=left] at (11.55,4.6) {$n_c=30$\\ crossover children};
  \node[font=\sffamily\scriptsize, anchor=west, align=left] at (11.55,0.8) {$n_m=8$\\ mutation children};
  \draw[decorate, decoration={brace, amplitude=5pt}, line width=0.7pt, gray!70!black] (13.9,7.55) -- (13.9,0)
       node[midway, right=7pt, font=\sffamily\scriptsize, align=left] {$P-n_e=38$\\ new evaluations\\ per generation};
\end{tikzpicture}%
}
\caption{Assembly of a new generation ($P=40$, defaults $\varepsilon=5\,\%$, $\chi=0.8$). The two
best designs survive unchanged; every other place is filled by a new child, bred from parents
chosen by binary tournament among \emph{all} $P$ individuals.}
\label{fig:split}
\end{figure}

\begin{table}[htbp]
\centering
\caption{Group sizes of Eq.~(\ref{eq:split}) for the default $\varepsilon=5\,\%$, $\chi=0.8$, and the
number of generations $G$ (Eq.~\ref{eq:generations}) for $B=2000$.}
\label{tab:split}
\begin{tabular}{@{}lcccccc@{}}
\toprule
 & $P$ & elites $n_e$ & crossover $n_c$ & mutation $n_m$ & evaluations / generation & $G$ \\
\midrule
$N_{\max}=5$  & 40 & 2 & 30 & 8  & 38 & 53 \\
$N_{\max}=10$ & 60 & 3 & 46 & 11 & 57 & 36 \\
\bottomrule
\end{tabular}
\end{table}

\paragraph{Elites are not the only parents.} Elitism decides which designs \emph{survive}
unchanged; selection decides which designs \emph{reproduce}, and every individual can be chosen.
In a binary tournament two individuals are drawn at random (with replacement) and the better one
wins, so the individual of rank $r$ (1 = best) wins a tournament with probability
\begin{equation}
\Pr(r)=\frac{2(P-r)+1}{P^2}.
\label{eq:tournament}
\end{equation}
With the $P-n_e=38$ tournaments of a generation of $P=40$, about 22 different individuals (56\,\%)
become parents. The best design is chosen 1.9 times on average, rank~10 1.5 times, rank~20 once,
rank~30 0.5 times and the worst almost never; three quarters of all parent choices fall on the
better half (Table~\ref{tab:tournament}). Good genes therefore spread through selection; the
elites only have to make sure that the best design itself is never lost.

\begin{table}[htbp]
\centering
\caption{Expected number of times an individual of rank $r$ is chosen as a parent in one
generation, $(P-n_e)\Pr(r)$ with Eq.~(\ref{eq:tournament}), for $P=40$, $n_e=2$.}
\label{tab:tournament}
\begin{tabular}{@{}lccccccc@{}}
\toprule
rank $r$ & 1 & 2 & 5 & 10 & 20 & 30 & 40 \\
\midrule
expected choices & 1.88 & 1.83 & 1.69 & 1.45 & 0.97 & 0.50 & 0.02 \\
\bottomrule
\end{tabular}
\end{table}

\paragraph{Why a small elite share.} One elite already guarantees that the best fitness never
gets worse, so the best-so-far curve of the \ac{GA} is monotone like that of \ac{SA}; a second
elite protects a second, different design, often one with another fleet size. More elites have
two costs: every elite place is one new design fewer per generation, and the elites, which
survive every generation \emph{and} are chosen as parents about twice each, quickly fill the
population with copies and near-copies of a few fleets (premature convergence). This matters here
because many designs are almost equally good: the forced-fleet-size study (validation V6) showed
that $J$ changes little between two, three and four spacecraft, and a small elite share keeps these
alternatives in competition.

\paragraph{Why separate crossover and mutation groups.} Each operator has one job. Crossover
\emph{combines} what the population already holds: a child can inherit a good polar orbit from one
parent and a good equatorial orbit from the other, a combination that \ac{SA} can only reach step
by step. Because whole slots are inherited, their simulations are already in the engine's cache and
a crossover child is cheap to evaluate. Slot crossover cannot, however, create an orbit radius or an
inclination that no parent has, and it changes the fleet size only by recombining switches;
mutation \emph{creates} new gene values and switches spacecraft on and off. The mutation children
make the same 1--2-gene move as \ac{SA}'s neighbour function, so they can be read as \ac{SA}-like
local steps and the crossover children as what the \ac{GA} adds on top. Keeping the groups separate
also lets every new best design be traced to the operator that produced it.

\paragraph{Parameters to be studied.} Both shares are treated as parameters and varied one at a
time at the same budget (Table~\ref{tab:split-study}). The elite share is expected to show a
U-shape: without elites the best design can be lost, with many elites the population converges
early. A crossover fraction of 1 removes the mutation group, so no new gene values can appear after
the first generation; a fraction of 0 leaves only mutation, which turns the \ac{GA} into many
\ac{SA}-like moves in parallel without a temperature.

\begin{table}[htbp]
\centering
\caption{Values of the elite share and the crossover fraction tested in the parameter study
(one parameter changed at a time; default in bold).}
\label{tab:split-study}
\begin{tabular}{@{}lccccc@{}}
\toprule
elite share $\varepsilon$ & 0 & 2.5\,\% & \textbf{5\,\%} & 10\,\% & 20\,\% \\
\quad elites $n_e$ for $P=40$ & 0 & 1 & \textbf{2} & 4 & 8 \\
\quad elites $n_e$ for $P=60$ & 0 & 2 & \textbf{3} & 6 & 12 \\
\quad generations $G$ for $P=40$ & 50 & 52 & \textbf{53} & 56 & 63 \\
\midrule
crossover fraction $\chi$ & 0 & 0.6 & \textbf{0.8} & 0.95 & 1.0 \\
\bottomrule
\end{tabular}
\end{table}

%=====================================================================================
\section{Crossover}
\label{sec:ga-crossover}

Crossover builds two children from two parents chosen by tournament. It works on whole slots,
in three steps: the slots of the second parent are \emph{aligned} to those of the first, each
aligned pair is then either \emph{swapped} or \emph{blended}, and the children are repaired
(Algorithm~\ref{alg:crossover}, Figure~\ref{fig:crossover}).

\subsection{Step 1: slot alignment}
The slots carry no identity: a fleet with orbits $(\mathcal{O}_1,\mathcal{O}_2,\mathcal{O}_3)$ in
slots 1--3 is the same mission as one with the orbits in slots 3, 1, 2. With $N_{\max}=5$ every
fleet can be written in up to $5!=120$ slot orders. Combining slot $i$ of one parent with slot $i$
of the other therefore usually pairs two \emph{unrelated} orbits: a child may receive the same
orbit twice (which then violates the separation constraint g3) or lose the orbits that made both
parents good. Before crossover, the slots of parent B are therefore reordered so that each slot of
parent A faces the most similar slot of B. The distance between two slots is
\begin{equation}
d(\mathbf{y}^A_i,\mathbf{y}^B_j)=\frac{|a^A_i-a^B_j|}{r_{\max}-r_{\min}}
+\frac{|\iota^A_i-\iota^B_j|}{180^\circ}
+\frac{\delta(\Omega^A_i,\Omega^B_j)}{180^\circ}
+\frac{\delta(M^A_{0,i},M^B_{0,j})}{180^\circ},
\qquad
\delta(\theta,\phi)=\min\bigl(|\theta-\phi|\bmod 360^\circ,\;360^\circ-|\theta-\phi|\bmod 360^\circ\bigr),
\label{eq:slotdist}
\end{equation}
where $\delta$ is the angle between two directions, so that $350^\circ$ and $10^\circ$ are
$20^\circ$ apart and not $340^\circ$. The pairs are matched greedily: flying slots of A with flying
slots of B first, closest pair first; then the remaining flying slots with switched-off ones;
finally the switched-off slots with each other. For $N_{\max}\le10$ this costs at most 100 distance
evaluations, nothing compared with one mission simulation, and needs no optimisation package.

\subsection{Step 2: swap or blend each aligned pair}
Each aligned pair of slots is treated in one of two ways:
\begin{itemize}
  \item \textbf{Swap} (probability $1-p_b$): child 1 takes the whole slot from one parent, chosen
        with probability $\tfrac12$, and child 2 from the other. The orbit is inherited intact,
        with its radius, inclination, node, phase and switch.
  \item \textbf{Blend} (probability $p_b$, only if \emph{both} slots fly): the orbit genes are
        interpolated with a weight $\lambda\sim U(-\alpha_b,\,1+\alpha_b)$ drawn per gene
        (blend crossover BLX-$\alpha$),
        \begin{equation}
        g^{(1)}=g^A+\lambda\,(g^B-g^A),\qquad g^{(2)}=g^B+\lambda\,(g^A-g^B),
        \qquad g\in\{a,\iota\},
        \label{eq:blx}
        \end{equation}
        and the angles along the shorter arc,
        \begin{equation}
        \Delta=\bigl((\theta^B-\theta^A+180^\circ)\bmod360^\circ\bigr)-180^\circ,\qquad
        \theta^{(1)}=(\theta^A+\lambda\Delta)\bmod360^\circ,\qquad \theta\in\{\Omega,M_0\}.
        \label{eq:arc}
        \end{equation}
        Both children keep the switch on.
\end{itemize}
The mission duration $T$ is treated like a slot: swapped with probability $1-p_b$, otherwise
blended with Eq.~(\ref{eq:blx}). A pair in which only one slot flies is always swapped, because an
average of an orbit that flies and one that does not has no meaning; the swap passes the flying
orbit to one child and the dormant one to the other, which is how crossover changes the fleet size.
Finally both children are passed through \texttt{repair} (reflection of $a$, $\iota$ and $T$ at
their bounds, which BLX-$\alpha$ can exceed, and wrapping of the angles), and a child identical to a
parent is mutated.

\paragraph{Why the angles need Eq.~(\ref{eq:arc}).} The plain average of $350^\circ$ and $10^\circ$
is $180^\circ$, the opposite direction, although the two nodes are only $20^\circ$ apart; the
shorter-arc rule gives $0^\circ$. Without it, two nearly identical orbits would produce a child on
the other side of the asteroid.

\paragraph{Why a random weight.} With the fixed weight $\lambda=\tfrac12$ (the midpoint of the two
parents) every blended gene lies strictly between its parents, and without selection the spread
of a gene in the population is halved each generation. After about five generations the
population would be confined to roughly $1/30$ of its initial spread, long before the end of the
run. With $\lambda\sim U(-\alpha_b,1+\alpha_b)$ a child can also land slightly beyond either
parent, which keeps the spread of the population; $\alpha_b=0.3$ is used.

\subsection{Choice of the blend probability}
Swapping and blending have different roles. Swapping \emph{recombines} complete orbits that have
already proved good: an orbit is good as a whole (radius, inclination and node together), and its
simulation is already in the engine's cache, so a swapped slot costs no new simulation. Blending
\emph{refines}: it creates an orbit between two matched good orbits, e.g.\ a radius of 57\,km between
parents at 55 and 60\,km, which swapping can never produce; each blended orbit is a new orbit to
simulate. Swapping should therefore dominate, with enough blending to keep refining the orbits.

A crossover child has about $N\approx3$ matched flying pairs (\ac{SA} chose two to four
spacecraft), so it contains on average $p_bN$ new orbits (Table~\ref{tab:pblend}). The default
$p_b=0.3\approx1/N$ gives about \emph{one} new orbit per child, which matches the mutation rate
(about one changed gene per child) and the cost of one \ac{SA} move (one changed orbit, one new
simulation), so that \ac{GA} and \ac{SA} also spend similar computing time per evaluation.
The blend probability is a parameter of the study, and the best value found there replaces this
default.

\begin{table}[htbp]
\centering
\caption{Effect of the blend probability $p_b$ on a crossover child with $N=3$ matched flying pairs.
The values marked $\star$ are tested in the parameter study, together with the fixed weight
$\lambda=\tfrac12$ and with the alignment switched off.}
\label{tab:pblend}
\begin{tabular}{@{}lcccc@{}}
\toprule
$p_b$ & $0^\star$ & $\mathbf{0.3}^\star$ & $0.5^\star$ & $1^\star$ \\
\midrule
new (blended) orbits per child, $p_bN$ & 0 & \textbf{0.9} & 1.5 & 3 \\
children without a new orbit, $(1-p_b)^N$ & 100\,\% & \textbf{34\,\%} & 12.5\,\% & 0\,\% \\
expected behaviour & \multicolumn{1}{>{\centering\arraybackslash}p{2.4cm}}{no refinement; only mutation improves an orbit}
 & \multicolumn{1}{>{\centering\arraybackslash}p{2.4cm}}{recombination and refinement}
 & \multicolumn{1}{>{\centering\arraybackslash}p{2.4cm}}{more refinement, more simulations}
 & \multicolumn{1}{>{\centering\arraybackslash}p{2.4cm}}{every orbit new; population contracts} \\
\bottomrule
\end{tabular}
\end{table}

\begin{algorithm}[htbp]
\caption{Slot crossover of two parents}
\label{alg:crossover}
\begin{algorithmic}[1]
\REQUIRE parents $\mathbf{x}^A,\mathbf{x}^B$; blend probability $p_b=0.3$; $\alpha_b=0.3$
\STATE \textbf{Align:} reorder the slots of $\mathbf{x}^B$ by greedy matching with Eq.~(\ref{eq:slotdist}) (flying--flying first)
\STATE $\mathbf{c}^{(1)}\leftarrow\mathbf{x}^A$; \quad $\mathbf{c}^{(2)}\leftarrow\mathbf{x}^B$ (aligned)
\FOR{each gene group $k\in\{T,\ \text{slot }1,\dots,\text{slot }N_{\max}\}$}
  \IF{$\mathrm{rand}()<p_b$ \textbf{and} group $k$ is $T$ or flies in both parents}
    \STATE \textbf{Blend:} $a,\iota,T$ with Eq.~(\ref{eq:blx}); $\Omega,M_0$ with Eq.~(\ref{eq:arc}); $\lambda\sim U(-\alpha_b,1+\alpha_b)$ per gene
  \ELSIF{$\mathrm{rand}()<\tfrac12$}
    \STATE \textbf{Swap:} exchange group $k$ between $\mathbf{c}^{(1)}$ and $\mathbf{c}^{(2)}$ (whole slot, switch included)
  \ENDIF
\ENDFOR
\STATE $\mathbf{c}^{(1)},\mathbf{c}^{(2)}\leftarrow\textsc{Repair}(\mathbf{c}^{(1)}),\textsc{Repair}(\mathbf{c}^{(2)})$
\STATE mutate any child that equals one of its parents
\RETURN $\mathbf{c}^{(1)},\mathbf{c}^{(2)}$
\end{algorithmic}
\end{algorithm}

\newcommand{\gaslot}[5]{% x, y, colour, line 1, line 2
  \filldraw[fill=#3!18, draw=#3!80!black, line width=0.8pt, rounded corners=3pt] (#1,#2) rectangle ++(2.3,0.8);
  \node[font=\sffamily\scriptsize, align=center] at ({#1+1.15},{#2+0.4}) {#4\\[-2pt]#5};}
\newcommand{\gaoff}[2]{% x, y
  \filldraw[fill=offc!10, draw=offc, dashed, line width=0.8pt, rounded corners=3pt] (#1,#2) rectangle ++(2.3,0.8);
  \node[font=\sffamily\scriptsize, text=gray] at ({#1+1.15},{#2+0.4}) {off (dormant)};}

\begin{figure}[htbp]
\centering
\resizebox{\textwidth}{!}{%
\begin{tikzpicture}[font=\sffamily, arr/.style={-{Stealth[length=5pt]}, line width=0.7pt}]
  % column x positions: 1.9, 4.5, 7.1, 9.7, 12.3
  \foreach \c in {1,...,5} { \node[font=\sffamily\scriptsize\bfseries, text=gray!60!black] at ({1.9+2.6*(\c-1)+1.15},7.55) {slot \c}; }

  % parent B as stored
  \node[font=\sffamily\small\bfseries, anchor=east, text=sc4!80!black] at (1.75,6.8) {Parent B (stored)};
  \gaslot{1.9}{6.4}{sc4}{41\,km\enspace$140^\circ$}{$\Omega=210^\circ$}
  \gaoff{4.5}{6.4}
  \gaslot{7.1}{6.4}{sc4}{52\,km\enspace$35^\circ$}{$\Omega=10^\circ$}
  \gaslot{9.7}{6.4}{sc4}{60\,km\enspace$90^\circ$}{$\Omega=130^\circ$}
  \gaoff{12.3}{6.4}

  % alignment: stored slot j -> aligned position i  (B3->1, B4->2, B1->3, B2->4, B5->5)
  \foreach \j/\i in {3/1, 4/2, 1/3, 2/4, 5/5} {
    \draw[arr, gray!70!black] ({1.9+2.6*(\j-1)+1.15},6.35) -- ({1.9+2.6*(\i-1)+1.15},5.25);
  }
  \node[font=\sffamily\scriptsize, fill=white, inner sep=1pt, text=gray!40!black, anchor=west] at (14.8,5.8) {\textbf{1. align}: nearest orbit, Eq.~(\ref{eq:slotdist})};

  % parent B aligned + parent A, stacked so matched pairs read as columns
  \node[font=\sffamily\small\bfseries, anchor=east, text=sc4!80!black] at (1.75,4.8) {Parent B (aligned)};
  \gaslot{1.9}{4.4}{sc4}{52\,km\enspace$35^\circ$}{$\Omega=10^\circ$}
  \gaslot{4.5}{4.4}{sc4}{60\,km\enspace$90^\circ$}{$\Omega=130^\circ$}
  \gaslot{7.1}{4.4}{sc4}{41\,km\enspace$140^\circ$}{$\Omega=210^\circ$}
  \gaoff{9.7}{4.4}
  \gaoff{12.3}{4.4}

  \node[font=\sffamily\small\bfseries, anchor=east, text=sc1!80!black] at (1.75,3.7) {Parent A};
  \gaslot{1.9}{3.3}{sc1}{50\,km\enspace$30^\circ$}{$\Omega=350^\circ$}
  \gaslot{4.5}{3.3}{sc1}{62\,km\enspace$95^\circ$}{$\Omega=120^\circ$}
  \gaslot{7.1}{3.3}{sc1}{40\,km\enspace$150^\circ$}{$\Omega=200^\circ$}
  \gaoff{9.7}{3.3}
  \gaoff{12.3}{3.3}

  % 2. swap or blend
  \foreach \c/\how in {1/blend $\lambda=\frac12$, 2/swap: from B, 3/swap: from A, 4/swap: from A, 5/swap: from B} {
    \draw[arr, gray!70!black] ({1.9+2.6*(\c-1)+1.15},3.25) -- ({1.9+2.6*(\c-1)+1.15},1.85)
        node[midway, fill=white, inner sep=1.5pt, font=\sffamily\scriptsize] {\how};
  }
  \node[font=\sffamily\scriptsize, fill=white, inner sep=1pt, text=gray!40!black, anchor=west] at (14.8,2.55) {\textbf{2. swap or blend} each pair};

  % child
  \node[font=\sffamily\small\bfseries, anchor=east] at (1.75,1.4) {Child};
  \gaslot{1.9}{1.0}{sc1!50!sc4}{51\,km\enspace$33^\circ$}{$\Omega=0^\circ$ (not $180^\circ$)}
  \gaslot{4.5}{1.0}{sc4}{60\,km\enspace$90^\circ$}{$\Omega=130^\circ$}
  \gaslot{7.1}{1.0}{sc1}{40\,km\enspace$150^\circ$}{$\Omega=200^\circ$}
  \gaoff{9.7}{1.0}
  \gaoff{12.3}{1.0}
  \node[font=\sffamily\scriptsize, anchor=west, align=left] at (14.75,1.4) {$N=3$; slots 2, 3 reuse\\cached simulations};
\end{tikzpicture}%
}
\caption{Slot crossover of two 3-spacecraft parents ($N_{\max}=5$; the duration $T$ is left out
for clarity). Parent B stores its orbits in a different slot order; without alignment, slot~1 of A
(50\,km, $30^\circ$) would be paired with slot~1 of B (41\,km, $140^\circ$), an unrelated orbit.
After alignment each column holds two similar orbits. Here the first pair is blended (the nodes
$350^\circ$ and $10^\circ$ give $0^\circ$ along the shorter arc) and the others are swapped. Blue:
from A; orange: from B; mixed: blended; grey: switched off.}
\label{fig:crossover}
\end{figure}

%=====================================================================================
\section{Mutation}
\label{sec:ga-mutation}

A mutation child is a copy of a parent in which one or two genes, chosen at random among the
\emph{active} genes (Section~\ref{sec:ga-chromosome}), are changed:
\begin{itemize}
  \item a real gene $g_j$ ($T$, $a_i$, $\iota_i$, $\Omega_i$ or $M_{0,i}$ of a flying spacecraft)
        receives a Gaussian step proportional to its range,
        \begin{equation}
        g_j\leftarrow g_j+\mathcal{N}\!\bigl(0,\;(\sigma_g\,(u_j-l_j))^2\bigr);
        \label{eq:mutation}
        \end{equation}
  \item a switch $s_i$ is flipped, which adds or removes a spacecraft. A slot that is switched on
        flies the dormant orbit it carries.
\end{itemize}
The child is then repaired. If repair maps it back onto the parent, or widens the change (switching
off the last spacecraft makes repair switch another one on), the mutation is drawn again. This is
the same move as the neighbour function of \ac{SA} (Section~\ref{sec:sa-problem}), so the mutation
children are \ac{SA}-like local steps.

The step size $\sigma_g$ used to breed generation $g+1$ from generation $g$ shrinks geometrically,
$g=0,\dots,G-2$,
\begin{equation}
\sigma_g=\sigma_0\Bigl(\frac{\sigma_{\min}}{\sigma_0}\Bigr)^{g/(G-2)},\qquad
\sigma_0=0.25,\quad \sigma_{\min}=0.01,
\label{eq:sigma}
\end{equation}
from a quarter of each variable's range (large jumps, exploration) to one percent (fine
adjustment). In \ac{SA} the step falls as $\sigma_0\sqrt{\Theta/\Theta_0}$ and the temperature falls
geometrically, so its step also shrinks geometrically from 0.25 to about 0.01; Eq.~(\ref{eq:sigma})
gives the \ac{GA} the same schedule, measured in generations instead of temperature levels.

\paragraph{Why only active genes.} A change to a dormant gene does not alter the fitness, so it
would cost an evaluation for nothing. Since every switch is an active gene, a switch is changed in
about $N_{\max}/(1+N_{\max}+4N)$ of the gene choices (5 of 18 for $N=3$, $N_{\max}=5$), so the
fleet size is mutated regularly.

\paragraph{Why one or two genes.} A mutation child should stay close to its parent, a good design,
so that the change can be judged: the effect of changing many genes at once is close to that of a
random design. Changing one or two genes matches \ac{SA}'s move and, with the crossover children
carrying about one new orbit each (Section~\ref{sec:ga-crossover}), keeps the cost of the
\ac{GA} per evaluation close to that of \ac{SA}.

%=====================================================================================
\section{Constraints, duplicates and termination}
\label{sec:ga-other}

\paragraph{Constraints.} Every child is evaluated through the same pipeline as in \ac{SA}: a
design that fails a pre-engine check (e.g.\ two spacecraft brought closer than $d_{\mathrm{sep}}$
by crossover) is not simulated and receives the penalised cost of Eq.~(\ref{eq:penalty}). Such a
child stays in the population: it loses every tournament against a feasible design, but among
infeasible designs the smaller violation wins, so the search is led back to the feasible region.
The separation and battery constraints cannot be repaired by moving genes into their bounds and are
handled by this penalty only.

\paragraph{Duplicates.} The elites are distinct designs, and a child identical to one of its
parents or to a child already created in the same generation is mutated again, so that no
evaluation is spent on a design whose fitness is known.

\paragraph{Termination.} The run stops when the evaluation budget $B$ is used, as for \ac{SA}; the
last generation may be incomplete. No stall criterion is used, so both algorithms always spend the
same number of evaluations. The result is the best design evaluated during the whole run.

%=====================================================================================
\section{One algorithm for the fleet and for benchmark functions}
\label{sec:ga-generic}

The \ac{GA} is validated, like \ac{SA}, on benchmark functions with a known optimum (validation V2).
These functions have no slots and no switches, so the problem-specific operators are attached to
an optional part of the \texttt{Problem} interface: the fleet pipeline reports which genes form each
slot (\texttt{slot\_groups}) and which genes are switches. A problem without this information is
treated with every gene as its own group: alignment is skipped, the slot swap becomes the standard
uniform crossover, the blend becomes the standard BLX-$\alpha$, and mutation changes one or two
genes. The same code therefore runs on both kinds of problem, and the benchmark results test the
\ac{GA} machinery (selection, elitism, crossover, mutation, budget) independently of the spacecraft
model.

%=====================================================================================
\section{The complete algorithm}
\label{sec:ga-algorithm}

Algorithm~\ref{alg:ga} and Figure~\ref{fig:gaflow} summarise the complete \ac{GA};
Table~\ref{tab:gaparams} lists its parameters.

\begin{algorithm}[htbp]
\caption{Genetic Algorithm for the fleet design}
\label{alg:ga}
\begin{algorithmic}[1]
\REQUIRE budget $B$; $P$ (Eq.~\ref{eq:popsize}); $\varepsilon$, $\chi$ (Eq.~\ref{eq:split}); $p_b$, $\alpha_b$; $k$; $\sigma_0$, $\sigma_{\min}$
\STATE \textbf{Start:} draw $P$ designs stratified over $N=1,\dots,N_{\max}$, each passing the pre-engine checks
\STATE evaluate them (Feasibility $\rightarrow$ Engine); $\mathbf{x}^{*}\leftarrow$ the best; $g\leftarrow0$
\STATE $n_e,n_c,n_m\leftarrow$ Eq.~(\ref{eq:split}); \quad $G\leftarrow$ Eq.~(\ref{eq:generations})
\WHILE{the budget is not used}
  \STATE $\sigma_g\leftarrow$ Eq.~(\ref{eq:sigma}); \quad $\mathcal{E}\leftarrow$ the $n_e$ best distinct designs \hfill (elitism)
  \STATE $\mathcal{C}\leftarrow\emptyset$
  \WHILE{$|\mathcal{C}|<n_c$}
    \STATE $\mathbf{x}^A,\mathbf{x}^B\leftarrow$ two different tournament winners
    \STATE add the children of \textsc{Crossover}$(\mathbf{x}^A,\mathbf{x}^B)$ (Algorithm~\ref{alg:crossover}) to $\mathcal{C}$
  \ENDWHILE
  \FOR{$n_m$ times}
    \STATE add \textsc{Repair}(\textsc{Mutate}(tournament winner, $\sigma_g$)) to $\mathcal{C}$ \hfill (1--2 active genes)
  \ENDFOR
  \STATE mutate again every child that duplicates a parent or another child
  \STATE evaluate the children until the budget is used; update $\mathbf{x}^{*}$ (feasible first, then lower $J$)
  \STATE population $\leftarrow\mathcal{E}\cup\mathcal{C}$; \quad $g\leftarrow g+1$
\ENDWHILE
\RETURN $\mathbf{x}^{*}$
\end{algorithmic}
\end{algorithm}

\begin{figure}[htbp]
\centering
\begin{tikzpicture}[
    font=\sffamily\small, node distance=0.45cm,
    box/.style={draw=#1, fill=#1!8, line width=0.9pt, rounded corners=4pt, text width=7.2cm,
                align=center, inner sep=5pt},
    box/.default=sc1,
    arr/.style={-{Stealth[length=6pt]}, line width=0.9pt, gray!70!black},
  ]
  \node[draw=headc, fill=headc, text=white, rounded corners=10pt, inner sep=6pt] (start) {Start};
  \node[box=headc, below=of start] (init) {\textbf{Initial population}: $P$ designs, stratified over the
        fleet size $N$, each passing the pre-engine checks};
  \node[box=headc, below=of init] (eval0) {\textbf{Evaluate} all $P$: Feasibility $\rightarrow$ Engine $\rightarrow J$};
  \node[diamond, aspect=2.6, draw=gray!70!black, fill=gray!8, line width=0.9pt, inner sep=1pt,
        below=of eval0] (dec) {budget left?};
  \node[box=headc, below=of dec] (elite) {\textbf{Elitism}: keep the $n_e$ best distinct designs};
  \node[box=sc1, below=of elite] (cross) {\textbf{Crossover} ($n_c$ children): tournament $\rightarrow$ 2 parents
        $\rightarrow$ align slots $\rightarrow$ swap or blend each pair};
  \node[box=sc4, below=of cross] (mut) {\textbf{Mutation} ($n_m$ children): tournament $\rightarrow$ 1 parent
        $\rightarrow$ change 1--2 active genes with step $\sigma_g$};
  \node[box=gray, below=of mut] (rep) {\textbf{Repair}, remove duplicates, \textbf{evaluate} the
        $P-n_e$ children; update the best design};
  \node[box=gray, below=of rep] (next) {next generation $=$ elites $+$ children; \ $g\leftarrow g+1$,
        shrink $\sigma_g$};
  \node[draw=headc, fill=headc, text=white, rounded corners=10pt, inner sep=6pt, align=center,
        right=1.6cm of dec] (stop) {return the best design\\ (feasible first, lowest $J$)};

  \draw[arr] (start) -- (init);
  \draw[arr] (init) -- (eval0);
  \draw[arr] (eval0) -- (dec);
  \draw[arr] (dec) -- node[right, font=\sffamily\scriptsize] {yes} (elite);
  \draw[arr] (elite) -- (cross);
  \draw[arr] (cross) -- (mut);
  \draw[arr] (mut) -- (rep);
  \draw[arr] (rep) -- (next);
  \draw[arr] (dec) -- node[above, font=\sffamily\scriptsize] {no} (stop);
  \draw[arr] (next.west) -- ++(-0.6,0) |- (dec.west);
\end{tikzpicture}
\caption{Flowchart of the Genetic Algorithm. Every evaluation passes through the same
Feasibility $\rightarrow$ Engine pipeline as \ac{SA}.}
\label{fig:gaflow}
\end{figure}

\begin{table}[htbp]
\centering
\caption{Parameters of the \ac{GA}: default values and the values compared in the parameter study
(one parameter changed at a time; the best value of each becomes the default).}
\label{tab:gaparams}
\footnotesize
\begin{tabularx}{\textwidth}{@{}Xccl@{}}
\toprule
Parameter & Symbol & Default & Values in the study \\
\midrule
Evaluation budget (per run) & $B$ & 2000 & --- (same as \ac{SA}) \\
Population size & $P$ & 40 / 60 (Eq.~\ref{eq:popsize}) & 10, 20, 40, 80, 120 \\
Initial population & --- & stratified & stratified, uniform \\
Elite share & $\varepsilon$ & 5\,\% & 0, 2.5, 5, 10, 20\,\% \\
Crossover fraction & $\chi$ & 0.8 & 0, 0.6, 0.8, 0.95, 1 \\
Tournament size & $k$ & 2 & 2, 3, 5 \\
Slot alignment & --- & on & on, off \\
Blend probability & $p_b$ & 0.3 & 0, 0.3, 0.5, 1 \\
Blend weight & $\lambda$ & $U(-\alpha_b,1+\alpha_b)$, $\alpha_b=0.3$ & random, fixed $\tfrac12$ \\
Genes changed per mutation & --- & 1 or 2 & --- (same as \ac{SA}) \\
Initial / final mutation step (fraction of range) & $\sigma_0$ / $\sigma_{\min}$ & 0.25 / 0.01 & $\sigma_0=0.1$, 0.25, 0.5 \\
Engine cache (orbits) & --- & 1024 & --- (does not change results) \\
\bottomrule
\end{tabularx}
\end{table}
